
\item \(f\) est le
morphisme structural \(\underline{\underline{P}}^r_Y\to Y\).
\end{enumerate}

\section*{4. Théorème de Riemann--Roch : cas d'une immersion fermée régulière}

\textbf{4.1.} Soit \(i:Y\to X\) une immersion fermée régulière. Le théorème de Riemann--Roch pour \(i\) est essentiellement la relation VII 4.4, qui permet de calculer les classes de Chern de \(i_*(x)\) à partir des \(\gamma^p(N,x)\). Plus précisément, si on introduit le \(\lambda\)-anneau \(K^{\bullet}(Y)_{\mathbb Q N}\) défini en V 5.5, alors l'homomorphisme de \(K^{\bullet}(Y)_{\mathbb Q N}\) dans \(K^{\bullet}(X)_{\mathbb Q}\) déduit de \(i_*\) est un \(\lambda\)-homomorphisme d'après VII 2.8 et VII 4.3. Le théorème de Riemann--Roch se ramène alors à une formule analogue concernant l'homomorphisme (additif) \(K^{\bullet}(Y)\longrightarrow K^{\bullet}(Y)_{\mathbb Q N}\) défini par \(y\rightsquigarrow(0,y)\) c'est-à-dire à une relation de nature formelle, valable pour tout \(\lambda\)-anneau. Nous ne développerons pas ici cette démonstration, afin d'éviter de faire encore du formalisme sur les \(\lambda\)-anneaux (cf. \(\RRR\) I 1.5), et nous reprendrons la méthode de [1] qui utilisait les schémas éclatés (méthode déjà employée pour la démonstration de VII 4.3, et qui est aussi celle de la fin de \(\RRR\)).

Comme il a été indiqué dans 3.9, on peut se limiter au cas d'une immersion régulière de codimension constante \(d\) vérifiant la condition de VII 4.1, le cas d'une immersion régulière générale se déduisant ensuite de 3.8. On reprendra dans ce paragraphe les notations de VII 3.1. On a donc le carré cartésien :
\[
\begin{tikzcd}[column sep=large,row sep=large]
Y' \arrow[r,"j"] \arrow[d,"g"'] & X' \arrow[d,"f"]\\
Y \arrow[r,"i"] & X ,
\end{tikzcd}
\]
